% Synced from content.md by sync_latex.py.
\section{\texorpdfstring{例 $1$ 四个条件，知三推一}{例 1 四个条件，知三推一}}

按图取定方位：$A$、$D$ 在 $BC$ 的同侧，$AC$ 与 $BD$ 在内部相交。延长 $CD$ 至 $X$，连接 $AD$，$DA$ 位于 $\angle BDX$ 内部。有四个条件：

\esdiagram{\esdiExterior}

\begin{center}
\small
\begin{tabularx}{\linewidth}{lL}
\toprule
条件 & 内容 \\
\midrule
\escond{1} & $AB=AC$ \\
\escond{2} & $\angle BAC=90^\circ$ \\
\escond{3} & $\angle BDC=90^\circ$ \\
\escond{4} & $DA$ 平分 $\angle BDX$，即 $\angle BDA=\angle ADX$ \\
\bottomrule
\end{tabularx}
\end{center}

\textbf{知道其中任意三个，就能推出剩下的一个。} 

\needspace{6\baselineskip}\subsection{\texorpdfstring{（$1$）$\text{\escond{1}\escond{2}\escond{3}}\Rightarrow\text{\escond{4}}$：连接直角顶点，得到外角平分线}{（1）①②③⇒④：连接直角顶点，得到外角平分线}}

已知 $AB=AC$，$\angle BAC=\angle BDC=90^\circ$，说明 $DA$ 平分 $\angle BDX$。

\think

两个直角先给出 $\angle ABD=\angle ACD$。接下来有两个方向：从 $A$ 向 $BD$、$CD$ 作双垂，把斜边 $AB$、$AC$ 配成直角三角形；或者围绕 $A$ 补出另一个等腰直角三角形，把上一章的手拉手用起来。

由 $\angle BAC=\angle BDC$，利用八字形倒角，得

\[
\angle ABD=\angle ACD.
\]

\needspace{6\baselineskip}\textbf{方法一：从 $A$ 作双垂。}

过 $A$ 作 $AM\perp BD$，垂足为 $M$；作 $AN\perp CD$，垂足为 $N$。本图中 $M$ 在线段 $BD$ 上，$N$ 在 $CD$ 经过 $D$ 的延长线上。

\esdiagram{\esdiDoublePerpendicular}

在 $\mathrm{Rt}\triangle ABM$ 与 $\mathrm{Rt}\triangle ACN$ 中，$AB=AC$，$\angle ABM=\angle ACN$，且两垂足处都是直角，故

\[
\triangle ABM\cong\triangle ACN\quad(\mathrm{AAS}),
\qquad AM=AN.
\]

再看 $\mathrm{Rt}\triangle ADM$ 与 $\mathrm{Rt}\triangle ADN$：斜边 $AD$ 公共，$AM=AN$，所以

\[
\triangle ADM\cong\triangle ADN\quad(\mathrm{HL}).
\]

因此 $\angle MDA=\angle NDA$，即 $\angle BDA=\angle ADX$。又因为 $\angle BDA+\angle ADX=90^\circ$，所以

\[
\boxed{\angle BDA=\angle ADX=45^\circ.}
\]

\needspace{6\baselineskip}\textbf{方法二：过 $A$ 添直角，构造手拉手。}

过 $A$ 作 $AP\perp AD$，交 $BD$ 于 $P$。

\esdiagram{\esdiHandLower}

由 $AB\perp AC$、$AP\perp AD$，按图有 $\angle BAP=\angle CAD$；又有 $\angle ABP=\angle ACD$、$AB=AC$，故

\[
\triangle ABP\cong\triangle ACD\quad(\mathrm{ASA}),
\qquad AP=AD.
\]

因此 $\triangle APD$ 为等腰直角三角形，$\angle PDA=45^\circ$。$P$ 在 $DB$ 上，故 $\angle BDA=45^\circ$，$\angle ADX=90^\circ-45^\circ=45^\circ$。

也可以先在 $BD$ 上取 $P$，使 $BP=CD$。由 $\angle ABP=\angle ACD$、$AB=AC$，先以 $\mathrm{SAS}$ 证明同一组全等，得到 $AP=AD$、$\angle BAP=\angle CAD$，再推出 $AP\perp AD$。两种作法落到同一幅图上：一种先添角，另一种先添边。

\insight{共斜边先倒角，再把“等腰”转成一组全等；外角平分线由等距离或等腰直角三角形得到。}

\needspace{6\baselineskip}\subsection{\texorpdfstring{（$2$）$\text{\escond{1}\escond{2}\escond{4}}\Rightarrow\text{\escond{3}}$：已知外角平分线，反推 $D$ 处直角}{（2）①②④⇒③：已知外角平分线，反推 D 处直角}}

已知 $AB=AC$，$\angle BAC=90^\circ$，$DA$ 平分 $\angle BDX$，说明 $\angle BDC=90^\circ$。

\think

这次不能先用两个直角倒角，因为 $D$ 处的直角正是待证结论。$\angle BDA=\angle ADX$ 能先制造对称：作双垂得到 $AM=AN$，或截取等长边制造 $\mathrm{SAS}$ 全等。把对称带来的等长再接上 $AB=AC$，便能推出 $\angle ABD=\angle ACD$。

\needspace{6\baselineskip}\textbf{方法一：作双垂，先用角平分线。}

沿用例 $1$（$1$）的 $AM$、$AN$。由 $\angle BDA=\angle ADX$、公共边 $AD$ 和两个直角，得

\[
\triangle ADM\cong\triangle ADN\quad(\mathrm{AAS}),
\qquad AM=AN.
\]

再结合 $AB=AC$，以 $\mathrm{HL}$ 得 $\mathrm{Rt}\triangle ABM\cong \mathrm{Rt}\triangle ACN$，所以 $\angle ABD=\angle ACD$。由八字形倒角，

\[
\angle BDC=\angle BAC=90^\circ.
\]

\needspace{6\baselineskip}\textbf{方法二：向上延长 $CD$，截取 $DE=DB$。}

延长 $CD$ 至 $E$，使 $DE=DB$，连接 $AE$。

\esdiagram{\esdiExtendCd}

由 $\angle BDA=\angle ADE$、公共边 $AD$、$DB=DE$，得

\[
\triangle ADB\cong\triangle ADE\quad(\mathrm{SAS}).
\]

所以 $AB=AE$，$\angle ABD=\angle AED$。又因为 $AB=AC$，所以 $AE=AC$，即 $\angle AEC=\angle ACE$。因此$\angle AEC=\angle ACE=\angle ABD$

最后八字形得 $\angle BDC=90^\circ$。

\needspace{6\baselineskip}\textbf{方法三：向右延长 $BD$，截取 $DF=DC$。}

延长 $BD$ 至 $F$，使 $DF=DC$，连接 $AF$。

\esdiagram{\esdiExtendBd}



\[
\angle ADF=180^\circ-\angle BDA,
\qquad \angle ADC=180^\circ-\angle ADX.
\]

等角的补角相等，所以 $\angle ADF=\angle ADC$。结合 $AD$ 公共、$DF=DC$，得

\[
\triangle ADF\cong\triangle ADC\quad(\mathrm{SAS}),
\qquad AF=AC,
\quad \angle AFD=\angle ACD.
\]

故 $AB=AC=AF$，从而 $\angle ABD=\angle AFD=\angle ACD$。最后八字形得 $\angle BDC=90^\circ$。

\insight{已知角平分线，可以先围绕它制造全等或者等腰。}

\needspace{6\baselineskip}\subsection{\texorpdfstring{（$3$）$\text{\escond{2}\escond{3}\escond{4}}\Rightarrow\text{\escond{1}}$：保留两个直角，反推等腰}{（3）②③④⇒①：保留两个直角，反推等腰}}

已知 $\angle BAC=\angle BDC=90^\circ$，$DA$ 平分 $\angle BDX$，说明 $AB=AC$。

\think

两个直角给 $\angle ABD=\angle ACD$；外角平分线给 $AM=AN$。这一次 $AB=AC$ 是待证结论，全等不能再用 $\mathrm{HL}$。已有两个角和一条直角边，可以改用 $\mathrm{AAS}$。

由 $\angle BAC=\angle BDC=90^\circ$，八字形倒角得 $\angle ABD=\angle ACD$。

从 $A$ 向 $BD$、$CD$ 作双垂，记垂足为 $M$、$N$。由 $\angle BDA=\angle ADX$，得 $\mathrm{Rt}\triangle ADM\cong \mathrm{Rt}\triangle ADN$，从而 $AM=AN$。

再由 $\angle ABM=\angle ACN$、两个直角、$AM=AN$，得

\[
\triangle ABM\cong\triangle ACN\quad(\mathrm{AAS}),
\qquad AB=AC.
\]

\textbf{另外两种对称全等的写法。}

\begin{itemize}
\item 按（$2$）方法二作 $E$，得到 $AB=AE$、$\angle ABD=\angle AED$。由 $\angle ABD=\angle ACD$，得 $\angle AEC=\angle ACE$，所以 $AE=AC$，进而 $AB=AC$。
\item 按（$2$）方法三作 $F$，得到 $AF=AC$、$\angle AFD=\angle ACD$。由 $\angle ABD=\angle ACD$，得 $\angle AFB=\angle ABF$，所以 $AB=AF$，进而 $AB=AC$。
\end{itemize}

\insight{条件与结论交换后，要重新选择全等判定；等边待证时，用等角和等距离搭桥。}

\needspace{6\baselineskip}\subsection{\texorpdfstring{（$4$）$\text{\escond{1}\escond{3}\escond{4}}\Rightarrow\text{\escond{2}}$：保留等边，反推 $A$ 处直角}{（4）①③④⇒②：保留等边，反推 A 处直角}}

已知 $AB=AC$，$\angle BDC=90^\circ$，$DA$ 平分 $\angle BDX$，说明 $\angle BAC=90^\circ$。

\think

$ABC$ 是否为直角三角形还不知道，不能提前使用手拉手或 $K$ 型全等。现在真正可用的是外角平分线和 $AB=AC$：先获得等距离，再用 $\mathrm{HL}$ 换出 $\angle ABD=\angle ACD$，最后倒角。

从 $A$ 向 $BD$、$CD$ 作双垂，垂足分别为 $M$、$N$。由 $\angle BDA=\angle ADX$、$AD$ 公共和两个直角，得 $\mathrm{Rt}\triangle ADM\cong \mathrm{Rt}\triangle ADN$，所以 $AM=AN$。

又 $AB=AC$，故

\[
\triangle ABM\cong\triangle ACN\quad(\mathrm{HL}),
\qquad \angle ABD=\angle ACD.
\]

由八字形倒角，得

\[
\angle BAC=\angle BDC=90^\circ.
\]

也可以照例 $1$（$2$）的两种延长作法：先由对称全等和 $AB=AC$ 得到 $\angle ABD=\angle ACD$，再把已知的 $\angle BDC=90^\circ$ 换回 $\angle BAC=90^\circ$。这两种作法在得出 $\angle ABD=\angle ACD$ 之前，都没有用到 $A$ 处是直角。

\insight{先清点已知条件；通过等距离与等斜边得到等角，再把直角搬回另一顶点。}

\needspace{6\baselineskip}\subsection{\texorpdfstring{单角变式：把\escond{4}换成 $45^\circ$ 或 $135^\circ$}{单角变式：把④换成 45° 或 135°}}

\needspace{8\baselineskip}\subsubsection*{\texorpdfstring{变式一：只给 $\angle BDA=45^\circ$}{变式一：只给 ∠BDA＝45°}}

已知 $AB=AC$，$\angle BAC=90^\circ$，$\angle BDA=45^\circ$，说明 $\angle BDC=90^\circ$。

\think

只给 $\angle BDA=45^\circ$，还没有给 $\angle ADX$。此时“到角两边距离相等”不能直接用。应让这个 $45^\circ$ 和新作的直角出现在同一个三角形里，先得到等腰直角三角形，再用手拉手；另一条路是把 $B$、$C$ 向 $AD$ 作双垂，用 $K$ 型全等搬运两条直角边。

\needspace{6\baselineskip}\textbf{方法一：过 $A$ 作 $AP\perp AD$，交 $BD$ 于 $P$。}

\esdiagram{\esdiHandLower}

由于 $\angle ADP=45^\circ$、$\angle PAD=90^\circ$，$\triangle APD$ 是等腰直角三角形，故 $AP=AD$。

由 $AB=AC$、$AP=AD$、$\angle BAP=\angle CAD$，得

\[
\triangle ABP\cong\triangle ACD\quad(\mathrm{SAS}).
\]

所以 $\angle ABP=\angle ACD$，即 $\angle ABD=\angle ACD$。由八字形倒角，$\angle BDC=\angle BAC=90^\circ$。

与例 $1$（$1$）对照：那里先全等，再证明 $AP=AD$；这里先由 $45^\circ$ 得 $AP=AD$，再全等。辅助线相同，证明顺序由已知条件决定。

\needspace{6\baselineskip}\textbf{方法二：从 $B$、$C$ 向 $AD$ 作双垂。}

过 $B$ 作 $BM\perp AD$，过 $C$ 作 $CN\perp AD$，垂足分别为 $M$、$N$。本构型中，沿 $AD$ 的方向，点的顺序为 $M$、$A$、$D$、$N$。

\esdiagram{\esdiKPerpendicular}

在 $\mathrm{Rt}\triangle ABM$ 与 $\mathrm{Rt}\triangle CAN$ 中，$\angle BAM=\angle ACN$，$AB=CA$，所以

\[
\triangle ABM\cong\triangle CAN\quad(\mathrm{AAS}),
\qquad BM=AN,\quad AM=CN.
\]

这里的角相等来自 $AB\perp AC$、$AM\perp CN$，按图取相等的锐角。

由于 $\angle BDM=45^\circ$，$\mathrm{Rt}\triangle BMD$ 为等腰直角三角形，所以 $BM=DM$。于是

\[
AN=DM=AM+AD,
\qquad DN=AN-AD=AM=CN.
\]

故 $\mathrm{Rt}\triangle CDN$ 也是等腰直角三角形，$\angle CDN=45^\circ$。$DN$ 与 $DA$ 反向，所以 $\angle ADX=45^\circ$。又 $\angle BDA=45^\circ$，故 $\angle BDC=180^\circ-45^\circ-45^\circ=90^\circ$。

\insight{只给单个 $45^\circ$ 时，先补等腰直角三角形；用 $K$ 型全等时，等长常藏在“加减公共部分”中。}

\needspace{8\baselineskip}\subsubsection*{\texorpdfstring{变式二：换成 $\angle ADC=135^\circ$}{变式二：换成 ∠ADC＝135°}}

已知 $AB=AC$，$\angle BAC=90^\circ$，$\angle ADC=135^\circ$，说明 $\angle BDC=90^\circ$。

\think

$\angle ADC=135^\circ$ 等价于 $\angle ADX=45^\circ$。与变式一对照，现在方便补出的是 $AD$ 与 $CD$ 延长线之间的等腰直角三角形。$K$ 型全等也仍可用，只是先得到等腰直角的变成了 $\triangle CDN$。

\needspace{6\baselineskip}\textbf{方法一：向上补等腰直角三角形。}

过 $A$ 作 $AQ\perp AD$，交 $CD$ 经过 $D$ 的延长线于 $Q$。

\esdiagram{\esdiHandUpper}

由 $\angle ADQ=180^\circ-135^\circ=45^\circ$，得 $AQ=AD$。又 $AB=AC$，且由两对垂直关系，按图有 $\angle BAD=\angle CAQ$，故

\[
\triangle ABD\cong\triangle ACQ\quad(\mathrm{SAS}).
\]

所以 $\angle ABD=\angle ACQ=\angle ACD$。再倒角得 $\angle BDC=90^\circ$。

\needspace{6\baselineskip}\textbf{方法二：从 $B$、$C$ 向 $AD$ 作双垂。}

沿用变式一的 $K$ 型图，由 $\mathrm{Rt}\triangle ABM\cong \mathrm{Rt}\triangle CAN$，有 $BM=AN$、$AM=CN$。

因为 $\angle CDN=180^\circ-\angle CDA=45^\circ$，所以 $\mathrm{Rt}\triangle CDN$ 为等腰直角三角形，$CN=DN$。于是

\[
AM=DN,
\qquad DM=AM+AD=DN+AD=AN=BM.
\]

故 $\mathrm{Rt}\triangle BMD$ 为等腰直角三角形，$\angle BDA=45^\circ$。已知 $\angle ADX=45^\circ$，所以 $\angle BDC=90^\circ$。

\insight{$135^\circ$ 先看邻补角；与 $45^\circ$ 变式比较，交换的是先出现的等腰直角三角形。}
